\documentclass[12pt]{article}

\usepackage[dvips,letterpaper,margin=0.75in,bottom=0.5in]{geometry}
\usepackage{cite}
\usepackage{slashed}
\usepackage{graphicx}
\usepackage{amsmath}
\usepackage{amssymb}
\usepackage{braket}
\usepackage[american,fulldiode]{circuitikz}

\begin{document}
\newcommand{\ihbar}{\ensuremath{i \hbar}}
\newcommand{\dPsidt}{\ensuremath{ \frac{\partial \Psi}{\partial t} }}
\newcommand{\dPsidx}{\ensuremath{ \frac{\partial \Psi}{\partial x} }}
\newcommand{\ddPsidx}{\ensuremath{ \frac{\partial^2 \Psi}{\partial x^2} }}
\newcommand{\dPssdt}{\ensuremath{ \frac{\partial \Psi^*}{\partial t} }}
\newcommand{\dPssdx}{\ensuremath{ \frac{\partial \Psi^*}{\partial x} }}
\newcommand{\ddPssdx}{\ensuremath{ \frac{\partial^2 \Psi^*}{\partial x^2} }}

\newcommand{\dphidt}{\ensuremath{ \frac{d \phi}{dt} }}
\newcommand{\dpsidx}{\ensuremath{ \frac{d \psi}{dx} }}
\newcommand{\ddpsidx}{\ensuremath{ \frac{d^2 \psi}{dx^2} }}


\date{\vspace{-5ex}}

\title{Homework Assignment 3 \\ Relativistic Kinematics}

\maketitle

\noindent
All problems are graded on effort only.\\

\noindent
{\bf Griffiths: P3.4, P3.9, P3.11, P3.15, P3.19, 3.21}\\[3pt]


\noindent
{\bf Problem 1:} Suppose a particle of mass $M$ has momentum $p$ and
energy $E$ in the lab frame.  It decays to three particles, A, B, and
C, each with mass $M/3$.  What is $\beta$ of particle A in the lab
frame?\\[3pt]

\noindent
{\bf Problem 2:}
Show that the process $\gamma \to e^- e^+$ cannot occur in the vacuum.\\[3pt]

\noindent
{\bf Problem 3:}
The anti-proton was discovered via the process:
$$p + p \to p + p + p + \bar{p}$$ Calculate the beam energy at which
this process is first kinematically allowed (A) for two colliding
beams of protons with equal energy, and (B) for one beam of protons
striking stationary protons.  Your calculation might lead you to
conclude that the antiproton discovery was made using colliding beams
(as in A), but in fact, the discover was made here in California,
using the bevatron in a fixed target experiment (as in B).  We did not
yet have the ability to focus beams which is necessary to collide two
beams at a high rate of collisions.\\[3pt]

\end{document}




